\pagebreak

\chapter{Using CPLDs}

    The Xilinx Coolrunner and XC9500 CPLDs are supported. The variable modes {\bf
    selectvalue}, {\bf queue}, {\bf cmemory} and {\bf rmemory} cannot be used.
    
    For {\bf static} mode variables the {\bf continuous} attribute is true by
    default. This means that top level static assignments will not have an
    initiated execution thread containing an infinite loop but instead the register
    clock enable will be continuously asserted.

    For {\bf output} mode variables the {\bf direct} attribute is true by default.
    This means that conditional (3-state) output will have the enable
    driven directly from the conditional test expression rather than from
    an initiated execution thread.
    
    Target multiplication, division, remainder and variable left or right shift
    are not implemented for CPLDs.
    
    The following 3PL program illustrates XC9500 usage -
    \begin{lstlisting}[gobble=4, language=threepl]
        part("XC9536PC44-5");

        module "xilinx/xc95.3pl";

        var(inpins, "[]str", {"P11", "P12", "P13", "P14", "P18", "P19", "P20"});
        str(outpin, "P22");
        var(countpins, "[]str", {"P34", "P35", "P36"});
        str(clkpin, "P5");

        clock(clk, frequency=1.0e6, loc=clkpin);

        useclock(clk);

        input(in, "[7]log", loc=inpins);

        static(o, "log");
        static(count, "uint:3", 5);
        output(,, count, loc=countpins);

        count++;

        seq {
            o <- in[0] && in[1] && in[2] && in[3] && in[4] && in[5] && in[6];
            o <- false;
        }

        output(,, o, loc=outpin);
    \end{lstlisting}
    
\pagebreak

    For Coolrunner CPLDs the I/O standard can be specified -   
    \begin{lstlisting}[gobble=4, language=threepl]
        part("XC2C32APC44-6");

        module "xilinx/xc2c.3pl";

        var(inpins, "[]str", {"P25", "P26", "P27", "P28", "P29", "P33", "P34"});
        var(outpin, "str", "P35");
        var(clkpin, "str", "P5");

        clock(clk, frequency=1.0e6, loc=clkpin);

        useclock(clk);

        input(in, "[7]log", loc=inpins);

        static(o, "log");

        seq {
            o <- in[0] && in[1] && in[2] && in[3] && in[4] && in[5] && in[6];
            o <- false;
        }

        output(,, o, loc=outpin);
    \end{lstlisting}
    
    The 3PL compiler will run the appropriate Xilinx CPLD fitting software,
    if post-processing is not disabled.
    The final product is a JEDEC file with a {\bf .jed} file name extension.
